\newcommand{\tipoRequerimento}{Contestação de Indeferimento por Desligamento de UC}

\section{DOS FATOS E DOS DIREITOS}

\begin{enumerate}[resume]
    \item Em resposta ao Protocolo nº 1815879064, referente à reclamação 002/2026 de Itaíba, a \nomeConcessionaria{} indeferiu e encerrou o pleito alegando, em suma, que \textit{“a instalação 6121379 foi desligada definitivamente a pedido do cliente. Desta forma a relação contratual foi encerrada para essa unidade consumidora”}.

    \item Contudo, a justificativa apresentada pela distribuidora para o indeferimento não possui qualquer amparo legal ou regulatório. O simples encerramento da relação contratual ou o desligamento definitivo da Unidade Consumidora (UC) \textbf{não extingue a responsabilidade da distribuidora} pelos débitos ou cobranças indevidas efetuadas durante o período de prestação do serviço.

    \item Se durante a vigência do contrato a UC 6121379 esteve enquadrada em classe/subclasse tarifária incorreta (em descumprimento ao disposto no art. 189 da REN nº 1.000/2021 da ANEEL, por se tratar de serviço público de iluminação pública/bens de uso comum do povo), restou caracterizado o faturamento a maior em prejuízo do erário público municipal.

    \item O dever de restituição decorre do princípio da vedação ao enriquecimento sem causa (art. 884 do Código Civil) e do regime de restituição de valores faturados a maior estabelecido na regulamentação do setor elétrico (art. 323 da REN nº 1.000/2021 da ANEEL). A extinção do contrato encerra a prestação continuada do fornecimento de energia, mas jamais o direito de regresso e ressarcimento do consumidor pelos valores indevidamente recolhidos no passado.

    \item Conforme destacado no item 35 e 36 da Nota Técnica nº 136/2024 da ANEEL, em casos de erro de faturamento imputável exclusivamente à distribuidora e constatados sob a vigência da REN nº 1.000/2021, resta impositiva a devolução dos valores pagos a maior, observando-se a regra da devolução em dobro quando caracterizada a ausência de engano justificável (§ 3º do art. 323 da REN nº 1.000/2021).

    \item O prazo prescricional aplicável para o ressarcimento de cobranças indevidas na prestação de serviço público de energia elétrica é de 10 (dez) anos, nos termos do art. 205 do Código Civil e ratificado pelo Despacho ANEEL nº 2.006, de 08 de julho de 2024. Portanto, o desligamento da UC não retira do Município o direito de obter a revisão dos faturamentos retroativos praticados no período decenal prescricional.

    \item Assim, revela-se totalmente genérica e improcedente a negativa da concessionária em analisar a revisão cadastral/faturamento retroativo com base unicamente na alteração do status da UC para “desligada”.
\end{enumerate}

\section{DOS PEDIDOS}

\begin{enumerate}[resume]
\item Ante o exposto, requer-se que a \nomeConcessionaria{}:
    \begin{enumerate}
        \item \textbf{Reconsidere e reforme o indeferimento} relativo ao Protocolo nº 1815879064, afastando o fundamento do desligamento definitivo da UC nº 6121379 como empecilho para análise do mérito do pedido de ressarcimento;
        \item Proceda à análise do correto enquadramento tarifário histórico da UC 6121379 (Classe Iluminação Pública - B4A), reconhecendo a cobrança indevida efetuada nos faturamentos passados;
        \item Efetue a restituição/devolução de todos os valores cobrados a maior referente aos últimos 10 (dez) anos de faturamento da referida UC, atualizados monetariamente pelo IPCA e acrescidos de juros de mora de 1\% ao mês calculados \textit{pro rata die}, conforme legislação e regulamentação da ANEEL;
        \item Promova a devolução em dobro das quantias cobradas indevidamente, nos termos do § 3º do art. 323 da REN nº 1.000/2021 e da Nota Técnica ANEEL nº 136/2024;
        \item Apresente o histórico das informações cadastrais da UC (data de ligação, datas e comprovantes de eventuais alterações tarifárias ocorridas ao longo do histórico de faturamento);
        \item Caso persista no indeferimento, apresente justificativa técnica e jurídica detalhada de mérito, acompanhada das devidas comprovações documentais (fotos, memória de cálculo, inspeções de carga instalada, etc.);
        \item Realize o pagamento/depósito do montante apurado diretamente em conta corrente de titularidade do Município de Itaíba-PE, conforme o § 7º do art. 323 da REN nº 1.000/2021, comunicando o resultado formalmente através do e-mail: \emailEmpresa{}.
    \end{enumerate}
\end{enumerate}


